% ==============================================================================
% Fichier : 01_cas_lizbiz.tex
% Description : Présentation du Cas Pratique Fil Rouge - Entreprise LiZbiZ
% ==============================================================================

\chapter{Cas Pratique Fil Rouge : Entreprise LiZbiZ}\label{chap:01_cas_lizbiz}

Ce chapitre présente le contexte de l'entreprise \textbf{LiZbiZ}, qui servira de support unique pour l'ensemble des travaux pratiques et mises en situation de ce cours.

\section{Identité et Contexte de l'Entreprise}

\subsection{Présentation Générale}

\textbf{LiZbiZ} est une entreprise commerciale spécialisée dans la vente de matériel informatique et high-tech. Elle s'est distinguée par une digitalisation rapide de ses processus, passant d'un modèle traditionnel à un modèle hybride fort sur le e-commerce.

\begin{itemize}
    \item \textbf{Secteur :} Commerce de gros et de détail / E-commerce.
    \item \textbf{Effectif :} 150 salariés.
    \item \textbf{Activité principale :} Vente de matériels, gestion de commandes en ligne, logistique associée.
\end{itemize}

\subsection{Contexte Stratégique et Enjeux}

L'entreprise traverse une phase critique de transformation.
\begin{enumerate}
    \item \textbf{Transformation Digitale :} LiZbiZ a récemment intégré des solutions d'Intelligence Artificielle (IA) pour optimiser sa logistique et son service client (Chatbots, prévision de stocks).
    \item \textbf{Exposition Médiatique :} L'entreprise est sous le feu des projecteurs. Des rumeurs de fraudes internes et de dysfonctionnements informatiques commencent à circuler.
    \item \textbf{Enjeu Réglementaire :} En tant qu'acteur majeur du e-commerce, LiZbiZ est soumise au RGPD et doit assurer la protection des données clients.
\end{enumerate}

\begin{boxlizbiz}
\textbf{Situation déclenchante :}
Suite à plusieurs articles de presse évoquant une possible fraude financière assistée par des failles informatiques et une utilisation contestable de l'IA, le Comité de Direction (COMEX) a décidé de mandat un audit approfondi pour rassurer les actionnaires.
\end{boxlizbiz}

% ------------------------------------------------------------------------------
\section{Le Système d'Information : Périmètre de l'Audit}
% ------------------------------------------------------------------------------

Le Système d'Information de LiZbiZ est centré autour d'une application majeure, simplement nommée \textbf{APPLICATION}, qui gère l'ensemble du cycle de vie de l'entreprise.

\subsection{L'Application "APPLICATION"}

C'est le cœur numérique de l'entreprise. Elle assure les fonctions vitales suivantes :
\begin{itemize}
    \item \textbf{Gestion Commerciale :} Création des comptes clients/fournisseurs, gestion des commandes.
    \item \textbf{Gestion Comptable :} Saisie des écritures, plan comptable, tenue des journaux.
    \item \textbf{Gestion Financière :} Facturation, suivi de la trésorerie, gestion des immobilisations.
    \item \textbf{Reporting :} Génération des états financiers et rapports de synthèse.
\end{itemize}

\subsection{Architecture Technique (Simplifiée)}

L'infrastructure repose sur une architecture cliente-serveur :
\begin{itemize}
    \item \textbf{Serveur de Base de Données :} SGBD Relationnel (PostgreSQL) hébergeant les données critiques.
    \item \textbf{Serveur d'Application :} Interface utilisateur et traitement métier.
    \item \textbf{Postes Clients :} Accès via navigateur web ou client lourd (comptabilité).
\end{itemize}

% ------------------------------------------------------------------------------
\section{Le Mandat d'Audit}
% ------------------------------------------------------------------------------

Cette section définit le cadre contractuel de votre mission en tant qu'auditeur.

\subsection{Origine de la Mission}

La mission fait suite à une demande formelle du COMEX, motivée par :
\begin{enumerate}
    \item La nécessité de vérifier la fiabilité des informations financières produites par l'APPLICATION.
    \item Le besoin de sécuriser le nouveau périmètre IA.
    \item La volonté de détecter d'éventuelles fraudes internes (Section 5.2 du cours).
\end{enumerate}

\subsection{Objectifs de la Mission}

L'objectif général est de réaliser un \textbf{Audit d'une Application en Service}. Les objectifs spécifiques sont :
\begin{itemize}
    \item \textbf{Fiabilité :} S'assurer que les données comptables sont exactes et complètes.
    \item \textbf{Sécurité :} Vérifier que les mécanismes de contrôle d'accès sont efficaces.
    \item \textbf{Conformité :} S'assurer du respect des procédures internes et des lois (RGPD).
\end{itemize}

\subsection{Organisation de l'Équipe (RACI)}

L'équipe de mission est composée comme suit :

\begin{table}[h!]
\centering
\begin{tabular}{lll}
    \toprule
    \textbf{Rôle} & \textbf{Responsabilité} & \textbf{Interlocuteur LiZbiZ} \\
    \midrule
    Chef de Mission & Supervision, Synthèse, Rapport & DSI / DG \\
    Auditeur Senior   & Tests techniques, Entretiens & Responsable Informatique \\
    Auditeur Junior   & Saisie, Extraction données   & Utilisateurs clés \\
    \bottomrule
\end{tabular}
\caption{Équipe d'audit et interlocuteurs}
\end{table}

% ------------------------------------------------------------------------------
\section{Environnement de Laboratoire (Immersion)}
% ------------------------------------------------------------------------------

Pour mener cette mission, nous disposerons d'un environnement simulé représentant le SI de LiZbiZ.

\subsection{Travaux Pratiques Préparatoires}

Dans le cadre de ce cours, vous serez amenés à :
\begin{enumerate}
    \item \textbf{Créer l'environnement :} Mettre en place la base de données de l'APPLICATION (Schéma relationnel fourni).
    \item \textbf{Charger les données :} Peupler la base avec un jeu de données réalistes (Clients, Factures, Écritures).
    \item \textbf{Concevoir les outils :} Développer des requêtes SQL et des feuilles de calcul pour automatiser les tests d'audit.
\end{enumerate}

\begin{boxalert}
\textbf{Mise en garde professionnelle :}
L'auditeur ne doit jamais intervenir sur le système de production réel sans précautions. L'utilisation d'un laboratoire ou d'une copie (snapshot) est impérative pour ne pas perturber l'activité de LiZbiZ.
\end{boxalert}

